% \numberwithin{equation}{section}

\chapter{Mètode} \label{sec:mètode}
En aquest apartat es detalla la implementació de les equacions en Dedalus i la resolució dels modes normals.
A més, es descriu també el tractament dels resultats.

Abans d'implementar els sistemes d'equacions obtinguts en la secció anterior %(Eq.~[\ref{eq.3.31}] per una banda i Eqs.~[\ref{eq.3.32}] i [\ref{eq.3.33}] per l'altra)
és indispensable adimensionalitzar les variables i, en conseqüència, les equacions.
D'aquesta manera es deixa de banda qualsevol sistema d'unitats i es proporcionen uns valors de referència adeqüats als valors típics de les variables del sistema:
\begin{align*}
	\bigl\{  l_{ref} \quad , \quad t_{ref} \quad , \quad v_{ref} \quad , \quad \rho_{ref} \quad , \quad B_{ref} \bigr\},
\end{align*}
on $v_{ref}$, $l_{ref}$ i $t_{ref}$ no poden ser escollides arbitràriament atès que han d'acomplir la condició $v_{ref} = l_{ref} / t_{ref}$.
Llavors les noves variables adimensionalitzades respecte els valors de referència són les següents:
\begin{equation}
	\begin{gathered}
		x = \bar{x} \cdot l_{ref}  \implies \frac{d}{dx} = \frac{1}{l_{ref}} \frac{d}{d \bar{x}} \quad , \quad \bar{\omega} = \omega \cdot t_{ref} \implies \omega = \frac{\bar{\omega}}{t_{ref}}, \\
		\rho = \bar{\rho} \cdot \rho_{ref} \quad , \quad v = \bar{v} \cdot v_{ref} \quad , \quad k_z = \frac{\bar{k_z}}{l_{ref}} \quad , \quad B = \bar{B} \cdot B_{ref}.
	\end{gathered}
	\tag{4.1} \label{eq.4.1}
\end{equation}

Finalment es substitueixen les variables normalitzades de l'Eq.~(\ref{eq.4.1}) dins l'Eq.~(\ref{eq.3.30}) en el cas del mode d'Alfvén,
i dins les Eqs.~(\ref{eq.3.31})~i~(\ref{eq.3.32}) en el cas dels modes ràpids i lents.
El resultat d'aquestes substitucions pràcticament no altera les equacions; únicament s'afegeix una barra sobre les variables.
Per qüestions visuals s'omet la barra i les equacions normalitzades que s'implementen al codi de Python són les Eqs.~(\ref{eq.3.30}), (\ref{eq.3.31}) i (\ref{eq.3.32}).

Es fa evident la subdivisió del problema degut al desacoblament de les equacions i, per tant, dels modes normals.
Aquests es descomposen en modes d'Alfvén i en ràpids i lents.

Per resoldre aquests dos subproblemes s'han creat diferents notebooks de Python (.ipynb) per separat en els que s'implementen les equacions corresponents,
apart del quadern de Mathematica. Tots els codis es poden trobar en l'Apèndix~\ref{ap.B}. %, veure \parencite{github}.

Per realitzar els càlculs es divideix l'espai del problema en $N = 256$ funcions base mitjançant polinomis de Legendre.
S'ha escollit aquest nombre per qüestions de memòria i temps de càlcul.

En el cas dels modes ràpids i lents, Dedalus resol les Eqs.~(\ref{eq.3.31}) i (\ref{eq.3.32}) per un valor de $k_z$ fixat.
Tant per aquests modes com pels d'Alfvén, Dedalus retorna els autovalors i les seves autofuncions.

Per això s'han de definir els límits, les constants i els camps del sistema.
Es defineixen les equacions dels problemes i les condicions de contorn i, finalment, es construeix un ``solver''.
Una vegada obtinguts els autovalors i les autofuncions, aquestes darreres es normalitzen al valor màxim del mode en valor absolut.
A partir d'aquest punt s'analitzen els resultats, en la Sec.~\ref{sec:resultats}.


\section{Modes d'Alfvén} \label{subsec:mètode_Alfvén}
El problema d'Alfvén de l'Eq.~(\ref{eq.3.30}) presenta una forma molt semblant al d'una corda amb els extrems fixats i amb estructura no uniforme.
La velocitat de propagació de les ones al llarg de la corda ve donada precisament per la velocitat d'Alfvén de l'Eq.~(\ref{eq.3.34}), que aplica a la Fig.~\ref{fig:sistema}.

Per aquest motiu i com a primer contacte amb la llibreria Dedalus \parencite{Dedalus},
s'ha resolt el problema dels modes normals d'una corda amb densitat no constant amb la condició de contorn d'extrems fixes.
Els resultats d'aquest exercici no s'afegeixen atès que són idèntics als del problema d'Alfvén,
que més envant es presenten en la~\hyperref[subsec:resultats_modes_Alfvén]{Subsec.~5.1}.


\section{Modes ràpids i lents} \label{subsec:mètode_r_i_l}
Una vegada resolt el primer problema es procedeix al càlcul de les autofuncions dels modes ràpids i lents, juntament amb la seva relació de dispersió.
En el cas d'acoblament nul amb velocitats purament longitudinal i vertical ($k_z = 0$) es recupera l'estructura del problema d'Alfvén,
com s'ha comentat en la~\hyperref[subsubsec:ones_modes_r_i_l]{Subsec.~3.2.2}.
És a dir, en aquest cas els modes ràpids i lents són idèntics als modes d'una corda en la que les ones es propaguen amb velocitat $v_A(x)$ i $c_s(x)$ de les Eqs.~(\ref{eq.3.34})~i~(\ref{eq.3.35}), respectivament.

%---------------------------------------------------------------------------------------------------------------------------------------------------
%\texttt{Prova codi}
%---------------------------------------------------------------------------------------------------------------------------------------------------

%---------------------------------------------------------------------------------------------------------------------------------------------------
%Comentari a fer amb ses etiquetes:
%per saber com son realment pots canviar sa cs o vA, o es tamany xp i també se veu més clar
%---------------------------------------------------------------------------------------------------------------------------------------------------